% Caption of figure F1_correctness_before_after of study_analysis_001 - generated by thesis/study_analysis/figures.py; do not edit by hand.
\caption[Correctness before and after repair (Q06, BASE vs FINAL)]{Correctness before and after repair (Q06, BASE vs FINAL). Population: BASE and FINAL[v] (identical base-cell cohort); unit: base cell; block: pilot, full\_ext, composite. Q06 correctness pass on BASE vs FINAL[v] (identical base-cell cohort, checked) per block. error bars: 95\% CI: nonparametric cluster bootstrap, percentile intervals, stratified by block (pilot benchmarks resampled among the 12 pilot benchmarks, full\_ext among its 48); seed 20261017; cluster = benchmark; 2,000 replicates (descriptive, unadjusted). BL01: Includes dense\_la/00 (BL-01 tolerance convention: possible false FAIL verdicts; disclosure-bearing aggregate); sensitivity SA1 in T10. IC1: Q06 is in-sample for test\_feedback and combined\_feedback: their loops see the ParEval test verdict and stop when it passes (every stopped\_tests\_pass trajectory is a Q06 pass), so they optimise the endpoint; for static\_feedback Q06 is out-of-sample (measured once in phase 2, never seen by the loop). The variants therefore differ in trigger and stop rules, not only in feedback content; test\_feedback cannot regress on Q06 by construction. IC6: static\_feedback repair triggers are dominated by clang-tidy bugprone findings and, on MPI, by low-confidence hints (design document: no trustworthy static MPI tool exists); static\_feedback results measure this frozen blocking-set configuration, not static analysis in general. \textsuperscript{a}~pilot: pilot\_002 / Pilotexperiment (12 benchmarks x 3 execution models x 11 models, 1 sample per prompt); pilot block has 12 clusters: percentile intervals are approximate; flagged in every pilot table.}\label{fig:study_analysis_001:F1_correctness_before_after}
